\documentclass{article}

\usepackage[swedish]{babel}
\usepackage{amsmath}
\usepackage{amssymb}

\title{Uppgift 2, Entydlighet av Maclaurin utveckling}
\author{Marcus Persson}
\date{\today}

\begin{document}

\maketitle

\section*{Uppgift}

Antag att funktionen $f$ och dess derivator
till och med ordning $n+1$ är kontinuerliga och att

\begin{equation}
\label{Maclaurin_utveckling}
    f(x) = a_0 + a_1x + a_2x^2
    + \cdots + a_nx^n + x^{n+1}B(x),
\end{equation}

\noindent där funktionen $B(x)$ är begränsad i en omgivning av $x=0$.
Visa då att:

\begin{equation}
\label{*}
    a_k = \frac{f^{(k)}(0)}{k!},
    \quad k = 0,1,2,\ldots,n.
\end{equation}


\section*{Bevis}

\noindent Enligt sats 9.2 kan $f(x)$ skrivas som

\begin{equation}
    f(x) =
    \sum_{k=0}^{n}
    \frac{f^{(k)}(0)}{k!}x^k
    + x^{n+1}B_1(x),
\end{equation}

\noindent där $B_1(x)$ är begränsad i en omgivning av $x=0$.

\noindent Av antagandet har vi även

\begin{equation}
    f(x) =
    \sum_{k=0}^{n} a_kx^k
    + x^{n+1}B(x),
\end{equation}

\noindent där $B(x)$ är begränsad i en omgivning av $x=0$.

\noindent Eftersom båda uttrycken beskriver samma funktion får vi

\begin{equation}
\begin{aligned}
    \sum_{k=0}^{n}
    \frac{f^{(k)}(0)}{k!}x^k
    + x^{n+1}B_1(x)
    &=
    \sum_{k=0}^{n} a_kx^k
    + x^{n+1}B(x).
\end{aligned}
\end{equation}

\noindent Vi använder induktions metoden för att visa
(\ref{*}) sant för $k=0,1,2,\ldots,n$.

\subsection*{Basfallet $k=0$:}

\begin{equation}
    a_0 = \frac{f^{(0)}(0)}{0!} = f(0).
\end{equation}

\noindent Antag att påståendet gäller för
$k=0,1,2,\ldots,p-1$, där $p$ fixed och $p\leq n$.

\subsection*{Visa uttalandet sant för $k=p$:}

\begin{equation}
\begin{gathered}
    a_px^p + \cdots + a_nx^n + x^{n+1}B(x)\\
    =\\
    \frac{f^{(p)}(0)}{p!}x^p
    + \cdots
    + \frac{f^{(n)}(0)}{n!}x^n
    + x^{n+1}B_1(x).
\end{gathered}
\end{equation}

\noindent Genom att dividera med $x^p$, då $x\neq 0$, får vi:

\begin{equation}
\begin{gathered}
    a_p + \cdots + a_nx^{n-p}
    + x^{n-p+1}B(x)\\
    =\\
    \frac{f^{(p)}(0)}{p!}
    + \cdots
    + \frac{f^{(n)}(0)}{n!}x^{n-p}
    + x^{n-p+1}B_1(x).
\end{gathered}
\end{equation}

\noindent När vi låter $x\longrightarrow 0$ får vi:

\begin{equation}
    a_p = \frac{f^{(p)}(0)}{p!}.
\end{equation}

\noindent Eftersom $B(x)$ och $B_1(x)$ är begränsade i en
omgivning av $x=0$ följer det av induktionsprincipen att:

\begin{equation}
    a_k = \frac{f^{(k)}(0)}{k!}.
\end{equation}

\noindent Där av följer enligt induktionsprincipen att
påståendet är sant $\forall\, k=0,1,2,\ldots,n$.
\begin{flushright}
    $\blacksquare$
\end{flushright}

\end{document}